% Transcription of folder 147, batch 08 (archivists' pages 141-160).
% Pass: Opus 5.5 (claude-opus-5-5), 2026-10-09 — first pass, unchecked against the pages by a human.
%
% Fast working hand throughout; the formulas carry better than the
% connective prose. Notation fixed for the batch: V_{D',Y} tubular
% neighbourhood, V^*_{D',Y} its trace on Y^* = Y \ \Theta, \Pi_1 fundamental
% groupoid, P_{D'^*} as written.
\documentclass[11pt,a4paper]{article}
\input{../preamble/grothendieck}

\folder{147}
\batch{8}
\pages{141}{160}
\foldertitle{Structure à l'infini des Mg,ν [pages 1 à 90, dont table des matières] : notes manuscrites (s.d.).}
\dating{s.d. — le groupe « Autour de La "Longue Marche" à travers la théorie de Galois » (141 à 149) est daté [à partir de 1978]-1983}
\watermark{Édition de démonstration}

\begin{document}

\page{141}
\note{p.~70 de l'auteur. La phrase commence au feuillet précédent, hors de ce lot.}
une définition directe, \struck{valable} dans un style encore plus topologique, valable chaque fois \ill{} \uncertain{(cf. ci-dessus)} — \ill{} une immersion locale de \emph{topos}. On va désigner par $V_{D',Y}$ le « voisinage tubulaire » défini \add{en fait} par une immersion \struck{top} locale topologique. Il a des propriétés de fonctorialité telles, que l'on trouve, \struck{dans} partant d'une situation \struck{de deux} d'un couple d'immersions locales $D''\to D'\to Y$, des diagrammes commutatifs
\[
(2.35)\qquad
\begin{array}{ccccc}
D'' & \longrightarrow & D' & & \\
\downarrow & & \downarrow & & \\
V_{D'',Y} & \xrightarrow[\text{imm. loc.}]{} & V_{D',Y} & \xrightarrow[\text{imm. loc.}]{} & Y
\end{array}
\]
\note{les flèches verticales $D''\to V_{D'',Y}$, $D'\to V_{D',Y}$ sont tracées comme des inclusions ; devant le diagramme, un renvoi biffé \struck{\ill{}}.}

\drawing{À gauche du diagramme, un croquis : une bande courbe hachurée transversalement, marquée $V_{D'}$, autour d'une courbe $D'$ ; deux points de $D''$ sur $D'$, entourés chacun d'un petit disque hachuré, désignés par des flèches pointillées comme $V_{D''}$. Un numéro biffé (\struck{(2.35)}) au-dessus.}

\struck{Dans} Pour l'instant on n'a pas utilisé le diviseur $\Theta$ à croisements normaux dans $Y$ (définissant l'ouvert $Y^*=Y\setminus\Theta$). Nous allons considérer maintenant l'image inverse de $Y^*$ dans $V_{D',Y}$ et $V_{D'',Y}$, soient $V^*_{D',Y}$ et $V^*_{D'',Y}$, s'insérant dans un diagramme

\page{142}
\[
(2.36)\qquad
\begin{array}{ccccc}
V^*_{D'',Y} & \longrightarrow & V^*_{D',Y} & \longrightarrow & Y\\
\| & & \| & & \\
V_{D'',Y}\setminus\Theta_{D'',Y} & & V_{D',Y}\setminus\Theta_{D',Y} & &
\end{array}
\]
\note{devant $V^*_{D'',Y}$, un début biffé \struck{\ill{}}. Sous $\Theta_{D'',Y}$, accolade et légende : « image inverse de $\Theta$ dans $V_{D'',Y}$ ».}

\marginal{\emph{NB} $V^*_{D'',Y}$ a la structure du complémentaire d'un diviseur à croisements normaux $\Theta''$, dans un voisinage tubulaire de \uncertain{$\Theta''(d_0)^*$} — et analogue pour $V^*_{D',Y}$.}

\drawing{Une bande courbe (le voisinage tubulaire), traversée par une courbe centrale ; en deux points, une branche transverse coupe la courbe, et autour de chaque point quatre quadrants hachurés ; une flèche désigne les branches comme $\Theta_{D',Y}$.}

Nous pouvons maintenant définir
\[
(2.37)\qquad \Pi_1 V^*_{D'',Y}\longrightarrow \Pi_1 V^*_{D',Y}
\]
en définissant les groupoïdes fondamentaux non pas par \struck{\uncertain{lacets}} \add{points et} / classes de chemins, mais par revêtements et isomorphismes entre tels\ldots

C'est essentiellement cela, l'\uncertain{heure} de canonisation — nous pouvons le définir chaque fois qu'on a deux des imm.\ locales $D''\to D'\to Y$

\page{143}
\note{p.~71 de l'auteur.}
(de topos simpliciaux, disons\ldots) et un ouvert $Y^*$ de $Y$ [plus généralement un objet du topos \struck{en question} $Y$, \uncertain{pas besoin} que ce soit un ss-objet de l'objet final, qui (dans le cas d'un tel ss-objet, i.e.\ définissant $Y^*$ comme ouvert du topos en question) \uncertain{à} supposer que $D'\to Y$ se factorise par $Y\setminus Y^*$]. Et on a des propriétés de transitivité \ill{} évidentes pour un composé
\[
(2.38)\qquad D'''\to D''\to D'\to Y .
\]

L'hypothèse « croisements normaux » va nous servir à expliciter les groupoïdes $\Pi_1 V^*_{D',Y}$, $\Pi_1 V^*_{D'',Y}$ en termes plus simples, comme \uncertain{canoniquement} équivalents à $\Pi_1 P_{D'^*}$, $\Pi_1 P_{D''^*}$. \emph{D}\ill{} il faut encore \add{\uncertain{(ici)}} expliciter la

\emph{Théorème.} Soient $Y$ un espace analytique, $\Theta$ un diviseur à croisements normaux de $Y$, $D'\to Y$ \struck{\ill{}} une immersion
\marginal{\ill{} — insertion écrite en oblique dans le coin inférieur gauche, rattachée à « Théorème » ; on y devine \uncertain{quitte à remplacer} \ldots\ $V_{D',Y}$ \ldots\ \uncertain{complexe} \ldots}

\page{144}
\marginal{\uncertain{Réduction de} plus haut \ill{} ici \ill{}}
\add{de codim $d$} (qui se factorise par $\Theta^{(d)}$ (\emph{NB} il est essentiel pour ce qui va \uncertain{suivre} \ill{} que $D'\subset\Theta^{(d)*}$), \struck{\ill{}} d'où un « voisinage tubulaire » $V_{D',Y}$ \struck{$*$} de $D'$ dans $Y$, et par trace de $Y^*=Y\setminus\Theta$ un \struck{\ill{}} germe d'espace analytique $V^*_{D',Y}$. Alors on a une équivalence de groupoïdes, \emph{définie à} \struck{équivalence} \emph{isomorphisme unique près}
\[
(2.39)\qquad \Pi_1 V^*_{D',Y}\xleftarrow{\ \approx\ }\Pi_1 P_{D'^*}
\]
\note{le parenthésage du début de page est celui du manuscrit ; l'exposant de $\Theta$ est d'abord écrit autrement puis surchargé en $(d)$.}

\emph{Remarques.} Le cas le plus simple est celui où $Y$ est de dim~1, et $d=1$, $D'$ réduit à un pt — \struck{le \ill{}} pas \uncertain{paraître} trivial, néanmoins la question de canonicité de l'équivalence m'a semblé nécessiter une \ill{}, et fait l'objet de la longue digression des « \uncertain{notes Monodromie} » §~53 n\textsuperscript{o}~5\textsuperscript{o} (p.~768' ff). Ici \add{(cas général)} il devrait être « assez évident » de définir \add{(au moins conceptuellement)} (un plongement « \emph{topologique} » \add{et via $C^\infty$} (non analytique!))

\page{145}
\note{p.~72 de l'auteur.}
\[
(2.40)\qquad V_{D',Y}\hookrightarrow \check N_{D',Y}
\]
\note{sous $\check N_{D',Y}$, accolade et légende : « fibré normal de $D'$ dans $Y$ ».}

\struck{transformant la \ill{}} \add{dont le composé avec} $Y\to V_{D',Y}$ soit la section nulle, et qui induise l'application identique sur les \struck{fib} fibrés normaux le long de la section nulle.\note{on lit bien $Y\to V_{D',Y}$ ; le sens demande sans doute $D'\to V_{D',Y}$.} \uncertain{Sa construction} \add{devrait} en tout cas dans le cas des vrais espaces analytiques paracompactes — c'est faux par contre dans le cas des multiplicités (Ex : \struck{$D'=E^1_{\mathbb C}\subset Y=E^2_{\mathbb C}$}, considérons $E^1\subset E^2$, avec l'opération du groupe discret $\mathbb Z$, \struck{plus} le générateur opérant par la transformation $\begin{pmatrix}1&1\\0&1\end{pmatrix}$, — pour $D'=(E^1;\mathbb Z)$, $Y=(E^2;\mathbb Z)$ — l'\struck{fibré} extension des fibrés vectoriels différentiables définie par le plongement $D'\subset V_{D',Y}\subset Y$ n'est pas scindable, donc pas celle définie par la section nulle d'un fibré l'est\ldots) Il est assez clair que lorsqu'on dispose d'une telle immersion \add{au moins pour des espaces analytiques ordinaires} (et sans condition de différentiabilité) celle-ci est une « équivalence d'homotopie », et
\note{les coefficients de la matrice sont en partie surchargés ; lecture $\begin{pmatrix}1&1\\0&1\end{pmatrix}$ probable. L'ajout « au moins pour des espaces analytiques ordinaires » est écrit dans la marge gauche, en face de « lorsqu'on ».}

\page{146}
induit donc une équivalence sur les $\Pi_1$. Mais le bât blesse, du fait qu'on n'a pas de canonicité — sans compter la nécessité d'hypothèses \add{restrictives} gênantes. \struck{Pour trouver} quelques doutes de canonicité \ill{}

\note{suit un diagramme, biffé par de grandes lignes ondulées : en haut $\Pi_1 P_{D'^*}\xrightarrow{\ \text{\uncertain{can}}\ }\Pi_1^{\text{rev}} V_{D',Y}$, en bas une flèche verticale marquée $\wr$ vers $\Pi_1^{\text{rev}} P_{D'^*}$, et à droite \emph{NB} $\Pi_1^{\text{\uncertain{rev}}}\approx$ \ill{}.}

\marginal{écrit en oblique dans la marge gauche, en partie surchargé : Mais pour définir $V^*_{D^*}\hookrightarrow V^*_{D'}$ \ill{} une équivalence entre les $\Pi_1$ de \ill{} et $P_{D'^*}$, il faut \ill{} \uncertain{rétablir} \ill{} les \uncertain{simplifications précédentes} \ill{}}

\ill{} \uncertain{d'entrer dans des} précisions d'ailleurs en relation avec le diviseur $\Theta$, quand nous \uncertain{prenons} le \struck{\ill{}} \add{plongement (2.40)} soit tel que l'image inverse de $P^*_{D'^*}$, identifié \add{($P_{D'}$, lui $=$ ouvert dans $\check N_{D',Y}$)} à un ouvert de $\check N_{D',Y}$, soit justement $V^*_{D',Y}$, de sorte qu'on \uncertain{trouve} \add{bien} \ill{}.
\[
\text{\struck{$\Pi_1 V^*_{D',Y}\longrightarrow \Pi_1 P_{D'^*}$}}
\]
\[
(2.41)\qquad V^*_{D',Y}\longrightarrow P_{D'^*}
\]
d'où
\[
(2.42)\qquad \Pi_1 V^*_{D',Y}\longrightarrow \Pi_1 P_{D'^*} .
\]
Pour avoir un morphisme de groupoïdes défini de façon plus « canonique », et

\page{147}
\note{p.~73 de l'auteur.}
sans hyp.\ superflues, je vais plutôt (comme dans Mon p.~779 ff) définir une flèche en sens inverse
\[
(2.43)\qquad \Pi_1^{\text{\uncertain{arcs}}} P_{D'^*}\longrightarrow \Pi_1^{\text{rev}} V^*_{D',Y}\ )
\]
où le premier $\Pi_1$ est défini par pts et arcs, et le deuxième (on n'a pas le choix!) est défini en termes de rev.~\uncertain{1-connexes}. Il faudrait des généralisations — \uncertain{adaptées} \ill{} les considérations de loc.\ cit.\ [je ne vais pas entrer dans le détail ici — d'autant moins que je ne vois pas \struck{clairement} à présent dans quelle mesure la th.\ \add{analogue} est indispensable — s'il n'est pas \uncertain{tout aussi} \add{bien commode} \struck{\ill{}} de travailler avec les $\Pi_1 V^*_{D',Y}$ — tirés des $\Pi_1 P_{D'^*}$ — en modifiant dans ce sens les définitions des $\widetilde\Pi_{d_0,d_1,\dots,d_r}$ ; Bien entendu, on aura besoin par la suite que $\widetilde\Pi_{d_0,\dots,d_r}\to\Pi_{d_0,\dots,d_r}$ ait une structure d'extension, de noyau un quotient de $T_{d_0,\dots,d_r}$ etc.,
\note{« Mon » abrège sans doute les notes sur la monodromie citées p.~144 (§~53) ; le renvoi est à la p.~779 de cette série.}

\page{148}
\note{le premier paragraphe est barré en entier de traits obliques ; il se lit néanmoins.}
\struck{Indépendamment du fait} \add{\struck{de \uncertain{sans} hyp.\ spécifique « crois.\ normaux »}} \struck{— on veut que $V^*_{D_1,D_0=X}$ puisse s'interpréter comme $V_{\Theta^{(1)},X}\setminus\Theta^{(1)}$ — il donne la « structure à l'infini ».}

Reprenons la situation du \uncertain{début} d'un pt de vue topologique (sans nous focaliser sur les $\Pi_1$). Le diviseur $\Theta$ à croisements normaux définit une filtration décroissante dans la multiplicité topologique $X$
\[
(2.44)\qquad \Theta^{(0)}=X\supset\Theta^{(1)}=\Theta\supset\Theta^{(2)}\supset\Theta^{(3)}\supset\cdots
\]
On peut définir
\[
D_d=\widetilde{\Theta^{(d)}}\quad\text{« normalisé topologique » de }\Theta^{(d)},
\]
dans un sens à préciser ----, \struck{\ill{}} mais cette définition dépend justement d'hyp « crois.\ normaux » \struck{et} relativement délicates. Par contre, \struck{\ill{}}
\note{la fin de la page, de « $D^*_d=$ » à (2.45), est barrée de traits obliques.}
\[
\text{\struck{$D^*_d=\Theta^{(d)}\setminus\Theta^{(d\pm1)}$}}
\]
\struck{se définit sans telles hyp\ldots, ainsi que}
\[
(2.45)\qquad \text{\struck{$V_{D^*_d}$ et $V^*_{D^*_d}\hookrightarrow V_{D^*_d}$}}
\]
\note{l'exposant du second $\Theta$ se lit mal : $(d+1)$ ou $(d-1)$ ; on écrit $d\pm1$ faute de mieux.}

\page{149}
\note{p.~74 de l'auteur. Le numéro (2.45) reprend celui de la formule barrée de la page précédente.}
Considérons le système ½ simplicial des
\[
(2.45)\qquad
\begin{array}{ccl}
 & & V^*_{d_0\cdots d_r}\overset{\text{déf}}{=}V^*_{D_{d_0\cdots d_r},X}\qquad\text{où }X=D_0\\
 & & \ \big\downarrow\\
D_{d_0,\dots,d_r} & \to & V_{d_0\cdots d_r}=V_{D_{d_0\cdots d_r},X}
\end{array}
\]
\marginal{où on a \ill{} $d_0\geq d_1\geq\cdots\geq d_r>0$}
\note{la flèche verticale $V^*_{d_0\cdots d_r}\to V_{d_0\cdots d_r}$ est tracée comme une inclusion.}

les $V^*_{d_0,\dots,d_r}$ s'envoient tous \struck{\ill{} dans} $X^*=D^*_0=$ \add{$X\setminus\Theta$} dans $V^*_0=X^*$, ils \struck{et sauf via $V^*_d\Rightarrow D$} s'envoient dans $V^*_{\Theta,X}$.

\emph{Th.\ de recollement.} Considérons le \struck{fibré en} \add{½ simplicial au-dessus de $V^*_{\Theta,X}$} topos \struck{(dont les composantes sont \ill{} plats)} \struck{(\ill{} un \ill{}} analytiques) formé des $V^*_{d_0,\dots,d_r}$ ($d_0\geq d_1\cdots\geq d_r>0$), \struck{Alors on} soit $V^*_{\bullet}$, on a un morphisme de la limite inductive des ces topos dans $V^*_{\Theta,X}$ \add{($=X^{*\infty}=D_0^{*\infty}$)}, et c'est là une équivalence de topos (en fait, c'est une équivalence d'homotopie, et en particulier se traduit en \uncertain{bijection} sur les $\pi_0$, et les $\pi_1$ relatifs à des \uncertain{rev.}~\ill{} qui se correspondent).

\page{150}
\emph{Corollaire 1.} Fixons un $\delta=\{\delta_0>\delta_1\cdots>\delta_s>0\}$, considérons $D_{\delta_0,\delta_1,\dots,\delta_s}=D_\delta$, \uncertain{le} diviseur $\Theta_\delta$ sur $D_\delta$, et $V^*_{\Theta_\delta,D_\delta}=D_\delta^{*\infty}$ (structure à l'infini de $D_\delta^{*\infty}$). Alors le syst.\ ½ simplicial des
\[
V^*_{D_{d_0,\dots,d_r,\delta_0,\delta_1,\dots,\delta_s}}\qquad(\text{pour }d_0>\cdots>d_r>\delta_0
\]
\note{sous les indices $d_0,\dots,d_r$, accolade et légende : « de degré \uncertain{simplicial} \uncertain{variable} » ; la parenthèse ouverte n'est pas refermée. Devant $V^*$, une lettre biffée \struck{\ill{}}.}
\uncertain{variables}) a comme \struck{\ill{}} \uncertain{lim} topologique la « structure à l'$\infty$ » $V^*_{\Theta_\delta,D_\delta}=D^*_\delta$.

\uncertain{Lemme} : \ill{} que le th.\ à $D_\delta,\Theta_\delta$\ldots

\note{la suite de la page, de « Cor.\ Main 2 » à la fin, est barrée de grands traits obliques ; elle se lit en partie.}
\struck{\emph{Cor.\ Main 2.} Supposons que}
\[
\text{\struck{$V^*_{\Theta,X}\hookrightarrow X^*$, \quad $V^*_{\Theta,X}\overset{\text{déf}}{=}X^{*\infty}$, \quad $X^*=X\setminus\Theta$}}
\]
\struck{soit une $\Delta$-équivalence (i.e.\ induise une bij.\ pour les $\pi_0$, et les $\pi_1$ des comp.\ conn.\ correspondantes) alors}
\[
\text{\struck{$\varinjlim\big(\widetilde\Pi_{d_0}\ (1\le d_0\le3),\ \widetilde\Pi_{d_0d_1}\ (d_0>d_1\geq$}}
\]
\note{la formule s'arrête au bas du feuillet ; devant $\varinjlim$, un premier membre raturé ($X^*$ \ill{}).}

\page{151}
\note{p.~75 de l'auteur.}
\emph{Corollaire 2.} On a
\[
(2.46)\qquad \Pi_1\,\text{\struck{$X$}}^{*\infty}\simeq\varinjlim\big(\widetilde\Pi_1\leftarrow\widetilde\Pi_{2,1}\rightarrow\widetilde\Pi_2\big)
\]
\note{sous $X^{*\infty}$, un signe d'égalité vertical renvoie à $V^*_{\Theta,X}$. Dans la limite inductive, $\widetilde\Pi_{2,1}$ est écrit au-dessus, ses deux flèches descendant vers $\widetilde\Pi_1$ et $\widetilde\Pi_2$.}

[En effet, une analyse locale facile (du type « th.\ de pureté ») montre qu'on peut \uncertain{dans un voisinage} négliger $\Theta^{(3)}$, i.e.\ que les \uncertain{voisins} pour \uncertain{nous} un \uncertain{changement} \ill{} si on remplace $X$ par $X\setminus\Theta^{(3)}$.

\marginal{écrit en oblique dans la marge gauche, en partie illisible : Il faut expliciter \ill{} \ill{} \uncertain{aussi} \ill{} le groupoïde $\widetilde\Pi_i$ ($i\geq3$) dans $\varinjlim$ \ill{} de $V^*_{\Theta,X}$ ; \uncertain{p.\ ex.}\ \ill{} le groupoïde $\widetilde\Pi_3$ est \ill{} $\varinjlim$ \ill{} \ill{} de type \ill{} \ill{} \ill{} groupoïde $\widetilde\Pi_{321}\to\widetilde\Pi_{3,2}$, $\widetilde\Pi_{3,1}$ \ill{} \ill{} $\varinjlim(\widetilde\Pi_1\leftarrow\widetilde\Pi_{21}\rightarrow\widetilde\Pi_2)$ \ill{}}

\emph{Corollaire 3.} Supposons $X^{*\infty}\to X^*$ soit une 1-équivalence, i.e.\ induise une équivalence pour les $\Pi_1$. Alors
\[
(2.47)\qquad \Pi_1 X\simeq\varinjlim\big(\widetilde\Pi_1\leftarrow\widetilde\Pi_{2,1}\rightarrow\widetilde\Pi_2\big).
\]
\note{sous la flèche $\widetilde\Pi_{2,1}\to\widetilde\Pi_1$, le mot « génération » \uncertain{(génération)} souligné.}

\note{la suite de la page est barrée de traits obliques.}
\struck{\emph{Corollaire 4.} Supposons que pour tout \add{$d\in\mathbb N$, \uncertain{et}} $i\leq d$, l'\uncertain{inclusion} soit}
\[
\text{\struck{$D_i^{*\infty}\hookrightarrow$ \ill{}, \quad $D_i^{*\infty}\overset{\text{déf}}{=}V^*_{\Theta_i,D_i}\setminus\Theta_i$}}
\]
\struck{Considérons une partie non multiplicité \add{$D_i$} d'un $D_i$ qui soit \uncertain{ouverte} et fermée, et considérons l'image $\Theta'^{(i)}\subset\Theta^{(i)}$ de $D_i$, qui est donc une « strate de \uncertain{codim} $i$ »,}

\page{152}
\marginal{en oblique, marge gauche : dire plutôt $\Sigma\subset X$ (sans exclure $\Sigma=X$) et demander qu'il ne \uncertain{passe} \uncertain{pas} \uncertain{exclusivement} \uncertain{les} \ill{} qui \ill{} \ill{} tel que $D_1^{\circ}$ \ill{} $D_0$}
Donnons nous une partie \add{\uncertain{fermée}} $\Sigma$ de $\Theta$, qui soit localement une réunion \struck{de de} \uncertain{de composantes} irréductibles des $\Theta^{(p)}$ ; une telle partie sera dite \emph{$\Theta$-admissible}. Si on désigne, pour tout $d\in\mathbb N^*$, par $D_d^\Sigma\subset D_d$ \struck{l'image inverse de $\Sigma$ dans $D_d$, par $D_d\to\Theta$, dans l'hyp.\ \ill{}} \add{la réunion des} \add{\uncertain{composantes}} \struck{\uncertain{irréductibles}} de $D_d$ qui s'envoient dans $\Sigma$, dans l'hyp.\ \uncertain{sur} $\Sigma$ \add{\uncertain{\ill{}}} \struck{\ill{}} : \struck{dire que} dire que $\Sigma$ est la réunion des
\[
(2.48)\qquad D_d^\Sigma\subset D_d
\]
\note{sous la formule, légende : « réunion de comp.\ irréd.\ de $D_d$ ».}
donc est connu quand on connaît les $D_\cdot^\Sigma$. Considérons un $d,d'$ ($d>d'\geq0$), d'où

\begin{tikzcd}[column sep=small, row sep=small]
 & D_{d,d'} \arrow[dl] \arrow[dr] & \\
D_d \arrow[dr] & & D_{d'} \arrow[dl] \\
 & D_0=X &
\end{tikzcd}

\note{la flèche $D_{d,d'}\to D_d$ porte une petite mention, lue \uncertain{étale}. À gauche du diagramme, un mot raturé \struck{\ill{}}.}
\[
D_{d,d'}^\Sigma\supset\text{l'image inverse de }D_{d'}^\Sigma,\qquad D_{d,d'}^\Sigma=\text{l'image inverse de }D_d^\Sigma
\]
\note{les deux membres sont disposés en colonne : « $D_{d,d'}^\Sigma\supset$ l'image inverse de $D_{d'}^\Sigma$ », et sous $D_{d,d'}^\Sigma$, un signe d'égalité vertical vers « l'image inverse de $D_d^\Sigma$ ».}

on voit donc que l'image inverse de $D_{d}^{\Sigma}$ dans $D_{d,d'}$, \struck{\ill{}}, est aussi égale à \struck{l'image} \struck{l'image dans} réunion des comp.\ \add{\uncertain{connexes}} \struck{irréductibles} de $D_{d,d'}$ qui s'envoient dans $\Sigma$, et \emph{contient} donc \uncertain{(mais n'est pas nécessairement égale)} l'image inverse de $D_{d'}^\Sigma$ dans $D_{d,d'}$. (Cela exprime que si une « branche » critique de codim $d'$ est dans $\Sigma$, les branches de codim \uncertain{$d$} \ill{} \uncertain{dedans} l'\uncertain{en} \uncertain{aussi}.)

\drawing{En bas à gauche, deux droites qui se coupent en un point marqué ; légendes $D_{2,1}$ : (flèche vers deux segments marqués d'un point), $D_2$ : un point, et une accolade $D_1$.}

\struck{La donnée de $\Sigma$} \add{d'une partie $\Theta$-admissible} \add{(de $\Theta$)} \struck{\ill{}} \struck{: \ill{} donnée \ill{}}

\page{153}
\note{p.~76 de l'auteur. Il n'y a pas de formule numérotée (2.49) sur ces feuillets.}
On \struck{désignera donc par $D_{d,d'}^\Sigma$ l'image} \add{posera donc}
\[
(2.50)\qquad D_{d,d'}^\Sigma=\text{image inverse de }D_{d'}^\Sigma\text{ par }D_{d,d'}\to D_{d'}
\]
(c'est \emph{pas} l'image inverse de $D_d^\Sigma$ par $D_{d,d'}\to D_d$ !), de sorte que $D_{d,d'}^\Sigma$ est aussi une réunion de comp.\ conn.\ de $D_{d,d'}$, et on a

\begin{tikzcd}[column sep=small, row sep=small]
 & D_{d,d'}^\Sigma \arrow[dl] \arrow[dr] & \\
D_d^\Sigma & & D_{d'}^\Sigma
\end{tikzcd}

\note{ce diagramme porte le numéro (2.51).}

La donnée d'une partie $\Theta$-admissible $\Sigma$ de $X=D_0$ est équivalente à la donnée d'un système de $D_d^\Sigma$, équivalente : la donnée de parties correspondantes
\[
(2.52)\qquad \pi_0(D_d^\Sigma)\overset{\text{déf}}{=}\pi_0(D_d)^\Sigma\subset\pi_0(D_d),
\]
satisfaisant la condition de compatibilité (2.51).
\note{le numéro est écrit « (2.51) » dans le texte ; devant (2.52), une rature \struck{$\pi_0$}.}

On désignera alors, pour tout drapeau $d_\bullet=\{d_0>d_1>\cdots>d_r\geq0\}$, par
\[
(2.53)\qquad D_{d_0\cdots d_r}^\Sigma=\text{image inverse de }D_{d_r}^\Sigma\text{ dans }D_{d_0,\dots,d_r}
\]
et on voit que \add{$D_{\Delta_r}^\Sigma=\coprod D_{d_0\cdots d_r}^\Sigma$} forment un système ½ simplicial de multiplicités analytiques, \struck{tel que} \uncertain{(qui correspond d'ailleurs à une sous-catégorie $\mathcal D_\bullet^\Sigma$ de la}
\note{l'ajout $D_{\Delta_r}^\Sigma=\coprod_{\ill{}}D^\Sigma_{d_0\cdots d_r}$ est écrit au-dessus de la ligne, l'indice de $\coprod$ illisible. La phrase continue au feuillet suivant.}

\page{154}
« catégorie analytique » \uncertain{envisagée} \uncertain{au début de ces} réflexions (p.~\uncertain{76} ff) — c'est \uncertain{en} fait une sous-catégorie pleine « stricte » (stricte signifiant que si un objet $y$ est, les objets qui s'envoient dedans y sont — c'est donc un \emph{crible}, en une terminologie \uncertain{vague} ?), dont la donnée équivaut à celle de \struck{\ill{}} $\Sigma$\ldots)

Ceci posé, le « th.\ de recollement » précédent se généralise en

\emph{Th.\ de recollement} (2\textsuperscript{ème} version). Les $\coprod V^{\Sigma\,*}_{D_{d_0\cdots d_r}}=V^{\Sigma\,*}_{\Delta_r}$ forment un syst.\ ½ simplicial de topos (de germes de multiplicités analytiques), dont la $\varinjlim$ est équivalente canoniquement à $V^*_{\Sigma,X}$.
\note{le chiffre des dizaines du renvoi se lit mal (76 ou 16) ; la p.~76 de l'auteur est la page 153 de ce lot.}

\section{12. Digression sur les déploiements des diviseurs à croisements normaux (suite)}

\page{155}
\note{p.~77 de l'auteur. Le titre est le sien ; il porte encore, biffé, \struck{\ill{} \uncertain{perspectives} canoniques associés à la \uncertain{stratification}, \ill{} \ill{} \ill{}} et se termine par « leur \uncertain{dissipation} en type \uncertain{élémentaire} ».}

\struck{\ill{}} Même dans le cas, \uncertain{relativement} \uncertain{ordinaire}, où $X=D_0$ est une courbe \add{\uncertain{entière}} alg.\ lisse (\struck{\ill{}} compacte p.~ex.), et où $\Theta$ se réduit : un diviseur lisse, le genre de formalisme développé dans la section précédente est remarquable. En fait, le « th.\ de recollement » est dans un pur \uncertain{tautologique}, \struck{\ill{}} \add{revenant} \struck{\uncertain{voisinage}} \add{\uncertain{essentiellement}} à dire que le voisinage tubulaire de $\Theta$ dans $X$, privé de $\Theta$, est la somme des \struck{voisinages} disques épointés correspondant aux pts de $\Theta$ ! Il y a cependant un \uncertain{intérêt} bien \uncertain{intéressant}, qui donne une \uncertain{interprétation} \struck{\ill{}} $\Pi_1(X)$ (et aussi $\Pi_1(X\setminus\Theta)$) comme somme \uncertain{amalgamée} de

\begin{tikzcd}[column sep=small, row sep=small]
 & \Pi_1(V^*_{\Theta,X}) \arrow[dl] \arrow[dr] & \\
\Pi_1 X^* & & \Pi_1 V_{\Theta,X}
\end{tikzcd}

\note{sous $\Pi_1 V_{\Theta,X}$, accolade et légende : « groupoïde \emph{discret} défini par $\Theta$ comme ens.\ de sommets ». Dans $\Pi_1(X\setminus\Theta)$, il avait d'abord écrit \struck{$X^*$} puis surchargé.}

i.e.\ de
\[
(2.54)\qquad \widetilde\Pi_0\longleftarrow\widetilde\Pi_{1,0}\longrightarrow\text{\struck{$\Pi_{1,0}$ \ill{}}}
\]
\note{la flèche de droite aboutit à un terme surchargé et biffé ; le diagramme continue au feuillet suivant.}

\page{156}
\note{les trois premières lignes sont barrées de traits obliques.}
\struck{où la flèche de droite n'a pas encore reçu droit de cité, dans le formalisme développé jusqu'à présent, des \ill{}} \struck{On voit en fait qu'il faut considérer},

En fait, dans le cas général, si
\[
(2.55)\qquad \Sigma\subset\Sigma'\ \text{\struck{\ill{}}}
\]
sont des inclusions de parties admissibles de $X$, on veut pouvoir exprimer $\Pi_1\Sigma$, $\Pi_1(\Sigma'\setminus\Sigma)$ \struck{\ill{}}, $\Pi_1V^*_{\Sigma,X}$, et \struck{$\Pi_1$ (\ill{} généralisation commune)} \add{plus généralement} \struck{\ill{}} $\Pi_1V^*_{\Sigma'\setminus\Sigma,X^*}$ en termes des $\widetilde\Pi_{\Delta_r}$ et des $\Pi_{\Delta_r}$ (et \uncertain{ainsi} introduire les éléments de structure qui manquent).
\note{dans $V^*_{\Sigma,X}$ et $V^*_{\Sigma'\setminus\Sigma,X^*}$, le $X$ est surchargé ; l'astérisque final se lit mal.}

Plus généralement, considérons
\[
(2.56)\qquad \text{\struck{\ill{}}}\ \Sigma'\subset\Sigma''\ \text{et un}\ \Sigma_1\subset\Sigma''
\]
(parties admissibles de $X$) et considérons $V_{\Sigma',\Sigma''}$, considérons la trace de $\Sigma_1$ dessus, qui peut s'écrire \struck{\ill{}}
\[
V_{\Sigma',\Sigma''}\cap\Sigma_1=V_{\Sigma'\cap\Sigma_1,\Sigma_1})
\]
\marginal{\emph{NB} on \uncertain{suppose} \ill{} que $\Sigma\subset\Sigma'$}
\note{l'indice de $\Sigma_1$ est surchargé dans les deux dernières occurrences ; « $\Sigma_1$ » est la lecture de la ligne (2.56).}

\page{157}
\note{le numéro de l'auteur en tête se lit « 75 », comme celui de la page 151 ; on attendrait 78.}
et considérons
\[
(2.57)\qquad V^{\Sigma_1}_{\Sigma',\Sigma''}\overset{\text{déf}}{=}V_{\Sigma',\Sigma''}\setminus V_{\Sigma',\Sigma''}\cap\Sigma
\]
\note{l'indice de $\Sigma_1$ en exposant est surchargé, et le dernier terme porte $\Sigma$ sans indice : sans doute $\Sigma_1$.}

Prenant \struck{\ill{}} $\Sigma'=\Sigma''$, on trouve $\Sigma'\setminus\Sigma$ (donc $\Sigma'$ si $\Sigma=\emptyset$) ; prenant $\Sigma''=X$, $\Sigma=\Theta$, on trouve $V^*_{\Sigma',X}$. On peut aussi considérer, en plus des données (2.56), un
\[
\Sigma'_0\subset\Sigma'
\]
et regarder
\[
(2.58)\qquad V^{\Sigma\setminus\Sigma\cap\Sigma'_0}_{\Sigma'\setminus\Sigma'_0,\ \Sigma''\setminus\Sigma'_0},
\]
\struck{mais} qu'on se gardera de confondre avec
\[
V^{\Sigma\cup\Sigma_0}_{\Sigma',\Sigma''}
\]
P.~ex.\ prenons $\Sigma=\emptyset$, $\Sigma'_0\subset\Sigma'\subset\Sigma''=X$, i.e.\ \uncertain{on a} $\Sigma'_0\subset\Sigma'\subset X$, il ne faut pas confondre
\[
V^{\Sigma_0}_{\Sigma',X}\quad\text{et}\quad V_{\Sigma'\setminus\Sigma_0,\ X\setminus\Sigma_0}
\]
p.~ex.\ si $\Sigma'_0=\Sigma'$, le premier terme est $V^{\Sigma'}_{\Sigma',X}$ (p.~ex.\ si $\Sigma'=\Theta$, c'est $V^*_{\Theta,X}$ \struck{\ill{}} qui pour un diviseur lisse est \struck{\ill{}} \add{équivalent} à un fibré en disques épointés sur $\Theta$) alors que le deuxième terme est le topos vide.

\page{158}
Notons que les topos du type (2.58) se ramènent au type (2.57), en remplaçant $X$, \struck{\ill{}} $\Sigma$, $\Sigma'$, $\Sigma''$, $\Sigma$ par \struck{\ill{}} $X\setminus\Sigma'_0$, $\Sigma'\setminus\Sigma_0$, $\Sigma''\setminus\Sigma'_0$, $\Sigma\setminus\Sigma\cap\Sigma'_0$. Pour faire cette réduction, il faut dans notre formalisme \uncertain{savoir} expliciter
\marginal{trait vertical double en face des deux lignes suivantes}
le passage au \emph{complémentaire} d'une partie admissible $\Sigma'_0$. Il faudra y revenir plus bas. D'autre part, les $V^{\Sigma}_{\Sigma',\Sigma''}$ se ramènent au cas $V^\Sigma_{\Sigma',X}$ (cas $\Sigma''=X$), pourvu qu'on sache expliciter le passage de $X$ à un $\Sigma''$ (bien entendu, on perd l'hyp.\ \uncertain{lisse} des croisements normaux — qui est un \uncertain{tic} \ill{}, — elle ne \uncertain{paraît} plus de rôle de \uncertain{puis} belle \uncertain{lurette} !\ldots) — c'est \uncertain{facile}, \uncertain{simplement}, par le système des $D^{\Sigma''}_{d_0\cdots d_r}$. \struck{\ill{}}

En fait, au niveau des topos, il

\page{159}
\note{le numéro de l'auteur se lit « 76 », comme celui de la page 153.}
suffit \add{en principe} de savoir exprimer les $V_{\Sigma',\Sigma''}$ et les morphismes de transition entre ceux-ci
\[
(2.59)\qquad V_{\Sigma',\Sigma''}\longrightarrow V_{\Sigma'_1,\Sigma''_1}\quad\text{pour}\quad(\Sigma',\Sigma'')\hookrightarrow(\Sigma'_1,\Sigma''_1),
\]
(quitte à appliquer ce formalisme, en remplaçant $X$ par $X\setminus\Sigma_0^{\ill{}}$ etc.) — on trouve dès lors $V^\Sigma_{\Sigma',\Sigma''}$ comme un \uncertain{sous-}topos \struck{\uncertain{résiduel}} \add{ouvert de $V^\Sigma_{\Sigma',\Sigma''}$} \uncertain{(relatif)} : l'image inverse de \struck{\ill{}} \uncertain{fermé} \uncertain{par} $V_{\Sigma'\cap\Sigma,\Sigma'}\longrightarrow V_{\Sigma',\Sigma''}$.
\note{la phrase est elliptique sur la page ; la place de « complémentaire » ou de « fermé » ne se lit pas sûrement.}
Mais cette « évidence » ne s'explicite pas aisément, au niveau des groupoïdes fondamentaux\ldots

Le fil directeur dans ce petit programme \uncertain{serait} celui-ci : on sait expliciter tous les topos qui sont associés à la filtration (comme ceux qu'on vient d'expliciter), et \uncertain{en} particulier leurs $\Pi_1$, en termes \struck{\ill{}} des topos $V_{d_0,\dots,d_r}\overset{\text{déf}}{=}V_{D^*_{d_0\cdots d_r},D_0}$ et $V^*_{d_0,\dots,d_r}=V^*_{D^*_{d_0\cdots d_r},D_0}$
\note{la ligne finale est surchargée ; les astérisques de $D^*$ et de $V^*$ sont ajoutés après coup, et la phrase continue au feuillet suivant.}

\page{160}
\marginal{en oblique, marge gauche : \ill{} \uncertain{suffit} \ill{} \uncertain{pour} \ill{}}
\uncertain{(Bien} \uncertain{entendu} \uncertain{la} \uncertain{notation} \uncertain{ne} \uncertain{fait} \uncertain{qu'il} s'agit de voisinages tubulaires des $D^*_{d_0,\dots,d_r}$, et non des $D_{d_0,\dots,d_r}$ eux-mêmes !) et des morphismes naturels de ces topos entre eux — quitte à y ajouter (\uncertain{par} \uncertain{la} \uncertain{même}) les topos analogues, obtenus en remplaçant
\[
\begin{array}{c}(X,\Theta)\\ \|\ \ \|\\ D_0\ \Theta_0\end{array}\quad\text{par un}\quad(D_{d_*},\Theta_{d_*})\ \text{quelconque}
\]
— ce qui conduit donc à envisager, en plus des $V_{d_0,\dots,d_r}$ et $V^*_{d_0,\dots,d_r}$, les topos analogues relatifs à une inclusion \struck{$d'_*=\{d'_0,\dots,d'_r$}
\[
d'_*=\{d'_0>\cdots>d'_r\geq0\}\hookrightarrow d_*=\{d_0>\cdots>d_r\geq0\}
\]
\[
(2.60)\qquad V_{d'_*,d'_*}\overset{\text{déf}}{=}V_{D^*_{d^*_*},\,D_{d'_*}}
\]
\note{sous le second membre, accolade et légende : « attention, on \uncertain{passe} au \uncertain{complémentaire} à $D^*_{d'_*}$ (mais non à $D_{d'_*}$ !) ». Dans les indices, des astérisques sont surchargés ; on transcrit ce que porte le dernier état. Devant (2.60) à gauche : « \uncertain{item} \uncertain{on} \uncertain{exprimant} $V^*$ \uncertain{aussi} ».}

qui pour $d'_*=d_*$ se réduit à $D^*_{d_*}$
\[
(2.61)\qquad V_{d_*,d_*}=D^*_{d_*}=D^*_{d_0,\dots,d_r},
\]
\uncertain{item} \ill{} et pour \struck{$d^*_*=\{d^*_0,\dots,d^*_r=0\}$ \uncertain{redondant} (\uncertain{quand} $d'_r\geq0$)} \add{\uncertain{redonnée}}
\[
V_{(d_1,\dots,d_{r-1},0),\,\Theta}=V_{(d_1,\dots,d_{r-1},0)}=V_{(d_1,\dots,d_{r-1})}
\]
\note{la dernière ligne est surchargée et serrée en bas du feuillet ; sous « item » on lit \uncertain{$V^*$ \uncertain{correspondant}}. L'argument se poursuit au-delà de ce lot.}
\end{document}
